% ============================================================================
%  The Bitcoin Block Space Problem:
%  Does the Fee Market Internalize Long-Term Storage Costs?
%  Working Paper v2.2.0 (model-spec.json v2.0.1) — LaTeX source
%
%  CONVERSION STATUS (honest):
%   - Abstract: FULLY CONVERTED from working-paper.md (verbatim).
%   - Sections 1-10: structure + all tables + key prose converted; some
%     long prose passages are condensed but faithful (see comments).
%   - References: converted from the paper's numbered list.
%   - NOT YET CONVERTED / STUBS: figures (none in the .md), the
%     verification-appendix content (kept as companion doc reference),
%     and full verbatim prose of every section (the .md remains the
%     canonical text). This source compiles; content parity with the .md
%     should be checked by diff before submission.
%
%  2026-08-03: subsection 7.1 (falsifiability) added to this source.
%  2026-08-03: fourth-reviewer deliverables applied --- assumption taxonomy
%  (Type C choices vs Type M mismeasurements) added as subsection 7.2; the v3.0
%  companion is labeled non-submission program material; advisor D5-first principle
%  recorded (roadmap §4); publication-plan v3.0 exclusion verified.
%  2026-08-03: peer-review execution applied --- abstract opening rewritten, SCCR
%  fraction equation + explicit-notation block (bundled C, time-dependent cb,
%  heterogeneous N, storage-vs-state) added to sec 4.1, average-vs-marginal
%  discussion in 4.2, 1x descriptive-not-normative stated early (sec 1/2/5.4),
%  voluntary-participation sentence in 8.1, limitations strengthened (sec 7).
%  2026-08-03: second-review restructure applied --- sections 11-12 of
%   working-paper.md were SPLIT into the companion research/future-directions-v3.md;
%   the core paper now ends at a compact section 10 (agenda table + pointer), which
%   is what this LaTeX source already mirrors. Abstract rewritten to range framing
%   (0.07-0.71 across N), fees-paid definition added (4.1), price-fee caveat (5.3),
%   current-N-band Monte Carlo added (5.4), terminology flipped to SCCR primary (8.3),
%   D5 status note added (abstract + conclusion).
%
%  TOOLCHAIN: pdflatex (not present on the dev machine at conversion
%  time — compile on any TeX installation:  pdflatex working-paper.tex)
%
%  AUTHOR: Prateek Poswal (identity choice pending — see
%  research/author-identity.md; swap the \author block once ratified).
% ============================================================================
\documentclass[11pt]{article}
\usepackage[margin=1in]{geometry}
\usepackage{amsmath,amssymb}
\usepackage{booktabs}
\usepackage{array}
\usepackage[T1]{fontenc}
\usepackage[utf8]{inputenc}
\usepackage[hidelinks]{hyperref}
\usepackage{microtype}
\usepackage{enumitem}

\title{\textbf{Storage Cost Internalization in Bitcoin's Fee Market}\\[0.4em] \large The Bitcoin Block Space Problem (program subtitle)}
\author{Prateek Poswal\\ \small Independent Researcher (Bitcoin Sahi Research)\\ \small ORCID: 0009-0005-2139-1877}
\date{Working Paper v2.2.0 \textbullet{} model-spec v2.0.1 \textbullet{} 2026-08-02}

\begin{document}
\maketitle
\vspace{-2em}
{\footnotesize This work is licensed under a Creative Commons Attribution 4.0 International License (CC BY 4.0): \texttt{https://creativecommons.org/licenses/by/4.0/}.}
\vspace{1em}

\begin{abstract}
Bitcoin has a market price for block space, but no explicit market price for long-lived resource consumption. This paper measures one of those resources --- replicated storage --- and asks how much of its modeled cost is covered by transaction fees. We define a \textbf{Storage Cost Coverage Ratio (SCCR)} --- the ratio of transaction fees paid (USD) to the estimated lifetime storage cost borne by full nodes (USD) --- and measure it against live fee-history data.

\textbf{The headline is a range, not a point.} Our best available measurement of the replication factor $N$ (full-node count) is a \textbf{primary-source lower-bound census ($\geq$32{,}000 known addresses via Bitcoin Core \texttt{getnodeaddresses})} --- the address-manager cap, not a complete enumeration --- while independent estimates span $\sim$10K--100K reachable nodes. Because SCCR is exactly inverse-linear in $N$, the result is stated in four lines (exact figures appear in \S\ref{sec:5}):

\begin{itemize}
\item \textbf{Dated baseline band (Aug 02--15): $\approx$0.22 coverage at $N=32$K}, with $\sim$100\% of sampled blocks below $1\times$ storage cost.
\item \textbf{Live series (Sep 07): 0.40 at $N=32$K, 94.37\% of blocks below $1\times$} --- the coverage ratio is \textbf{fee-market-driven and time-varying}, not a fixed point.
\item \textbf{Dynamic finding:} SCCR moved \textbf{0.16 $\rightarrow$ 0.45 $\rightarrow$ 0.40} across a \textbf{22-point daily series (2026-08-02 $\rightarrow$ 2026-09-07)} measured live in \texttt{data/sccr\_history.json} --- dynamic measurement strengthens the paper.
\item \textbf{Model uncertainty: depends strongly on the replication factor $N$} --- the average spans $\sim$0.71 at $N=10$K to $\sim$0.07 at $N=100$K; this $N$-band range ($\sim$0.07--0.71) is a \emph{different} uncertainty from the dated/live observations above and carries the dominant risk.
\end{itemize}

A joint Monte Carlo over $N$, $C$, $T$, and price (current band, 10{,}000 draws) gives a P5--P95 interval of 0.07--0.47, median 0.17, with 99.9\% of draws below the $1\times$ threshold; the share of sampled blocks below $1\times$ ranges $\sim$79--100\% depending on the true $N$. We reconcile two cost models that previously disagreed by $16.4\times$ (dimensionless), document the correction transparently (model-spec.json v2.0.0), and state results to the precision the evidence licenses: \textbf{external reproduction is pending (D5)}. The framework is reproducible, falsifiable research.
\end{abstract}

\noindent\textbf{Hypothesis:} Bitcoin's fee market efficiently allocates scarce block space, but may not fully internalize every long-lived resource cost created by confirmed transactions.

\section{Introduction: Congestion Pricing vs.\ Permanence Cost}\label{sec:1}
The one-sentence frame (full treatment in \texttt{problem\_statement.md}):
\begin{quote}
\textbf{Bitcoin's fee market prices competition for inclusion in the next block, but it does not explicitly price the long-term resource costs of permanently recorded blockchain data.}
\end{quote}
The fee market solves a short-term optimization problem extremely well: during congestion, higher fees win block space. This is \textbf{congestion pricing}, with a horizon of roughly one block ($\sim$10~min). Storage cost is structurally different:
\begin{itemize}[nosep]
  \item \textbf{One-time payment} (the fee, USD/block) vs.\ \textbf{recurring cost} (long-term replicated storage across the network, USD/(node$\cdot$yr))
  \item \textbf{Payer chooses} to pay vs.\ \textbf{node operators bear} the cost involuntarily
  \item $\sim$10-min horizon (next block) vs.\ \textbf{indefinite horizon} (data persists in blockchain history)
\end{itemize}
\textbf{Scope:} this paper measures the \textbf{storage leg} of a broader resource-pricing question. Bandwidth, validation, and UTXO-maintenance legs now carry v1 bounds/measurements (\S\ref{sec:5} 5.5--5.7) with the residual program in \S\ref{sec:7}.

\textbf{Why storage? It is not the ``most important'' resource --- it is simply the first measurable one.} Storage qualified first because its cost leg ($C$, $N$, $T$) and its fee leg (block fees, USD) are both estimable from primary sources. Every other resource is a named research hypothesis, not a measured result.

\textbf{Storage $\neq$ state.} SCCR measures the cost of replicated \emph{history} --- the permanent record of confirmed blocks retained by full nodes. It does \emph{not} measure the \textbf{UTXO set}, the live ledger state every node maintains in RAM and index structures. State permanence is a separate resource with its own accounting leg (UCIR, \S\ref{sec:7} / roadmap Phase II).

\textbf{$1\times$ is a descriptive calibration point, not a normative target.} Stated at the outset so no reader mistakes the paper for a claim that fees \emph{should} cover 100\% of modeled storage cost. The paper measures whether they do; it prescribes nothing.

\section{What We're NOT Saying}\label{sec:2}
\begin{itemize}[nosep]
  \item $\times$ Not proposing a fix
  \item $\times$ Not claiming the externality is economically significant at current volumes (it may not be)
  \item $\times$ Not claiming the SegWit discount was a mistake (it solved transaction malleability)
  \item $\times$ \textbf{Not arguing that Bitcoin is `broken'} --- the fee market solves block-space allocation well; whether it internalizes long-lived resource costs is an empirical question
  \item $\times$ \textbf{Not claiming fees \emph{should} cover 100\% of storage cost} --- the $1\times$ threshold is a descriptive calibration point, not a normative target (\S\ref{sec:1}, \S\ref{sec:5.4})
  \item $\checkmark$ Just framing the question clearly and measurably
\end{itemize}

\section{The SegWit Pricing Structure}\label{sec:3}
SegWit (BIP~141, activated August 2017) introduced block weight: witness data counts \textbf{1 weight unit (WU) per byte}, non-witness data \textbf{4~WU per byte}. This \textbf{4:1 (dimensionless) discount} reduces the marginal block-space cost of witness-resident data relative to non-witness data. Years later, inscription protocols (Ordinals, BRC-20, Runes) took advantage of that pricing structure. Two precise caveats (full treatment in \texttt{bip141\_analysis.md}):
\begin{enumerate}[nosep]
  \item The discount applies to \textbf{all} witness data --- SegWit financial transactions benefit identically. The ``accident'' is that it also subsidizes data-bearing constructions, not that SegWit was designed for them.
  \item Whether the discount appropriately reflects the long-term resource costs of the data it makes cheaper is \textbf{part of the research question}, not a foregone conclusion.
\end{enumerate}

\section{The Model}\label{sec:4}
\subsection{Canonical specification}\label{sec:4.1}
All quantities, units, and formulas live in \texttt{research/model-spec.json} (v2.0.1) --- a single canonical source that every script imports. No script redefines a model constant.
\begin{center}
\renewcommand{\arraystretch}{1.25}
\begin{tabular}{lllll}
\toprule
Quantity & Symbol & Units & Value & Source \\
\midrule
Annual node cost & $C$ & USD/yr & 925 & \texttt{utxo\_cost\_model} component sum (924.35) rounded \\
Replication factor & $N$ & nodes & 32{,}000 & primary-source lower-bound census ($\geq$32{,}000 known addresses via Bitcoin Core \texttt{getnodeaddresses}; addrman-saturated --- \textbf{a lower bound}; independent estimates 10K--100K) \\
Storage horizon & $T$ & yr & 10 & assumption (archival retention) \\
Average block size & $B_{\mathrm{block}}$ & bytes & 1{,}500{,}000 & \texttt{block\_stats} / \texttt{fee\_history} \\
Blocks per year & $R_{\mathrm{blocks}}$ & blocks/yr & 52{,}596 & $365.25 \times 24 \times 6$ \\
Total block bytes/yr & $B_{\mathrm{all,yr}}$ & bytes/yr & $7.8894\times10^{10}$ & $B_{\mathrm{block}} \times R_{\mathrm{blocks}}$ \\
\textbf{Cost per byte per year} & \textbf{$c_b$} & USD/(byte$\cdot$yr) & \textbf{$1.17246\times10^{-8}$} & $C / B_{\mathrm{all,yr}}$ (block-average) \\
Lifetime storage cost/node/block & $L$ & USD/block & 0.175869 & $c_b \times B_{\mathrm{block}} \times T$ \\
Network lifetime cost/block & $L_{\mathrm{net}}$ & USD/block & 5{,}628 & $L \times N$ ($=32$K) \\
\textbf{Storage Cost Coverage Ratio} & \textbf{SCCR} & dimensionless & 0.2228 (167 blocks, live 2026-08-02); $\sim$0.29 (156 blocks, dated 2026-08-01); $\sim$0.07--0.71 across the N=10K--100K band & $\mathrm{fee}_{\mathrm{USD}} / L_{\mathrm{net}}$ \\
\bottomrule
\end{tabular}
\end{center}
\begin{center}
\renewcommand{\arraystretch}{1.5}
\begin{tabular}{c}
\rule{0pt}{14pt}$\displaystyle \mathrm{SCCR} = \frac{\text{Fees Paid (USD/block)}}{\text{Modeled Lifetime Storage Cost (USD/block)}} = \frac{\mathrm{fee}_{\mathrm{USD}}}{L_{\mathrm{net}}} = \frac{\mathrm{fee}_{\mathrm{USD}}}{C \cdot T \cdot N / R_{\mathrm{blocks}}}$\rule{0pt}{14pt}
\end{tabular}
\end{center}

\textbf{Definition of ``fees paid'' (used once, everywhere).} Throughout this paper, \emph{fees paid} means \textbf{total transaction fees of the block in USD}: $\mathrm{fee_{USD}}(block) = \mathrm{avgFees}(block) \times \mathrm{USD/BTC\ price}$, where $\mathrm{avgFees}$ is the block's total transaction fees in satoshis from the live \texttt{fee\_history} capture (mempool.space 24h block-fee history, into \texttt{captured-data/bsahi.db}). It \textbf{excludes the block subsidy} (miner revenue, not a fee) and is \textbf{not} a fee \emph{rate} (sat/vB). All three independent implementations (JS, Python, standalone C) compute $\mathrm{fee_{USD}} = \mathrm{avgFees}/10^8 \times \mathrm{USD}$ identically. The denominator it is compared against is $L_{\mathrm{net}} = C \times T \times N / R_{\mathrm{blocks}}$ (USD/block).

\textbf{Design principle:} $c_b$ is \textbf{horizon-free} ($C /$ annual bytes). The horizon $T$ enters \textbf{only} through $L = c_b \times B \times T$. This corrects the v1.0.0 implementation, which applied $T$ twice (see \S\ref{sec:6}).

\textbf{Explicit notation (hidden assumptions made visible).} The canonical quantities bundle four assumptions the model treats as scalars; each is stated explicitly so readers see exactly what is and is not modeled. None changes the arithmetic of \S\ref{sec:5}:

\begin{enumerate}[nosep]
  \item \textbf{$C$ is a bundled cost, not pure storage.} $C = C_{\mathrm{storage}} + C_{\mathrm{bandwidth}} + C_{\mathrm{misc}}$ --- the \$925/yr annual node cost bundles disk/SSD, bandwidth, and electricity/hardware depreciation (component sum 166.67 + 600 + 157.68). The model treats $C$ as a single bundled scalar, so the headline ratio is strictly a \textbf{storage-and-hosting coverage ratio}. The bandwidth-vs-storage decomposition is a documented future refinement --- exactly the RCIR and BCIR legs (\S\ref{sec:7} / roadmap Phase II--III).
  \item \textbf{$c_b$ is time-dependent, not stationary.} The canonical $c_b$ is a point-in-time ratio $c_b(t) = C(t)/B_{\mathrm{year}}(t)$. $B_{\mathrm{year}}(t)$ depends on block fullness: 95\%-full blocks produce $\sim$7.5$\times$10$^{10}$ bytes/yr vs $\sim$5.5$\times$10$^{10}$ at 70\%-full blocks, moving $c_b$ (and SCCR) even with $C$, $N$, $T$ fixed. SCCR as measured is therefore a \textbf{point-in-time estimate over its capture window} (\S\ref{sec:5.1}).
  \item \textbf{$N$ is a heterogeneous vector, not a scalar.} The model writes $L_{\mathrm{net}} = L \cdot N$; the explicit forward notation is $L_{\mathrm{net}} = \sum_i L_i$, where each node class $i$ (archival, pruned, exchange, research/indexer) carries different storage behavior and hence a different per-node lifetime cost $L_i$. This is precisely the archival-vs-pruned measurement gap documented in the companion note \texttt{archival-vs-pruned-note.md}: the census measures reachable count $N$ ($\geq$32K, a lower bound), not the retention distribution.
  \item \textbf{Storage $\neq$ state (UTXO).} SCCR prices replicated \emph{history} --- the permanent record of confirmed blocks --- not the UTXO set, the live ledger \emph{state} nodes maintain in RAM and index structures. State permanence is a separate resource with its own accounting leg (UCIR, \S\ref{sec:7} / roadmap Phase II).
\end{enumerate}

\subsection{The marginal-inscription attribution (secondary)}\label{sec:4.2}
The inscription-externality branch uses a \textbf{marginal attribution}: $c_{b,\mathrm{insc}} = C / (\text{inscription bytes/yr}) = 1.92573\times10^{-6}$~USD/(byte$\cdot$yr). This differs from $c_b$ by a factor of \textbf{164.4$\times$ (dimensionless)} --- not an error, but two different denominators (all block bytes vs.\ inscription-only bytes). The paper headline uses the block-average $c_b$; the marginal figure appears only in the inscription-externality analysis. The $16.4\times$ (dimensionless) figure reported in earlier versions was this $164\times$ (dimensionless) denominator gap \textbf{divided by} the $10\times$ (dimensionless) bug --- see \S\ref{sec:6}.

\textbf{Why average, not marginal?} $c_b$ is an \emph{average} (accounting) attribution: total node cost divided by total bytes. The $164\times$ (dimensionless) marginal branch ($c_{b,\mathrm{insc}}$) attributes the same cost to inscription bytes only --- a \emph{marginal} (optimization) view. Both are valid; they answer different questions. Average cost is the right tool for \textbf{accounting} --- ``what share of the network's modeled hosting bill does the fee market cover?'' --- which is this paper's question. Marginal cost is the right tool for \textbf{optimization} --- ``what does one more byte cost the network?'' --- which is the inscription-externality branch's question. The paper uses the average for its headline (accounting) and documents the marginal variant explicitly, rather than burying the choice.

\section{Findings}\label{sec:5}
\subsection{The Storage Cost Coverage Ratio}\label{sec:5.1}
Measured from live \texttt{fee\_history} captures at node count $N=32{,}000$ --- the best-available \textbf{primary-source lower-bound census} ($\geq$32{,}000 known addresses via Bitcoin Core \texttt{getnodeaddresses}; \S\ref{sec:5.4}). The headline, in three lines (exact numbers in the tables below):

\begin{itemize}
\item \textbf{Representative live measurement: $\approx$0.22}
\item \textbf{Observed band: $\approx$0.22--0.29} (across captures at $N=32$K)
\item \textbf{Model uncertainty: depends strongly on the replication factor $N$} --- the true-$N$ band ($\sim$0.07--0.71, \S\ref{sec:5.4}) is a \emph{different} uncertainty from the observed-sample band above
\end{itemize}

Two snapshots, dated explicitly:
\begin{center}
\renewcommand{\arraystretch}{1.25}
\begin{tabular}{lcc}
\toprule
Metric & Dated capture (2026-08-01) & Live capture (2026-08-02) \\
\midrule
Blocks sampled & 156 & 167 \\
Average SCCR (dimensionless) & $\sim$0.293 & 0.2228 \\
Min / Max (dimensionless) & $\sim$0.058 / $\sim$1.537 & $\sim$0.058 / $\sim$0.832 \\
Blocks below $1\times$ & 98.7\% (154/156) & 100.0\% \\
\bottomrule
\end{tabular}
\end{center}
The ratio moves with the fee market; the single source of truth is \texttt{research/model-spec.json} (v2.0.1, canonical); all surfaces must read the live value from \texttt{node tools/research/storage-ratio.js}, never a hardcoded figure. The v2.0.0 N=60K-era values (0.1719 / 0.1535) are superseded.

\textbf{Fee-regime + node-count dependence:} the ratio tracks both the fee market and the replication factor. Under the earlier N=60K assumption the average was 0.156--0.172 (dimensionless) with 100\% below $1\times$; at the primary-source lower-bound census N=32K it rises to $\sim$0.22--0.29 with $\sim$98.7--100\% below $1\times$ (a few high-fee blocks exceed coverage in the dated capture). The headline is a \emph{distribution over time and parameters}, not a point.

\subsection{The inscription externality (marginal branch)}\label{sec:5.2}
At 100{,}000 inscriptions/month ($\sim$400 bytes each, in witness):
\begin{center}
\begin{tabular}{lc}
\toprule
Metric & Value \\
\midrule
Cost per byte per year (marginal) & $1.92573\times10^{-6}$ USD/(byte$\cdot$yr) \\
Lifetime storage cost per inscription (10 yr) & \$0.00770 USD/inscription \\
New 10-yr storage liabilities created per year at 100K/mo & $\sim$\$9{,}240 USD/yr spread across all nodes \\
Steady-state annual externality (amortized) & $\sim$\$924 USD/yr spread across all nodes \\
\bottomrule
\end{tabular}
\end{center}
\textbf{Unit note (corrected):} \$9{,}240 is the \emph{undiscounted 10-year liability of one year's inscriptions} (USD/yr), not an annual cost. Amortized over the 10-yr horizon it is $\sim$\$924/yr. This was previously labeled ``annual unpriced externality'' --- a T/1 conflation of the same shape as the fixed $10\times$ bug, in label only.

\textbf{Pruning-consistency note:} this externality is structurally real but \textbf{not economically significant at current volumes} --- per-node it is well under \$1/yr. The contribution of this paper is the reproducible ratio framework, not the magnitude.

\subsection{Sensitivity}\label{sec:5.3}
Recomputed at the live capture (2026-08-02, 167 blocks; baseline SCCR $=$ 0.2228 (dimensionless) at N=32K, C=\$925/yr, T=10~yr). The table's parameter ranges bracket the canonical values:
\begin{center}
\renewcommand{\arraystretch}{1.2}
\begin{tabular}{lcl}
\toprule
Parameter & Value & Avg SCCR (dimensionless) \\
\midrule
Node cost, $C$ (USD/yr) & 600 / 925 / 1400 & 0.343 / 0.223 / 0.147 \\
Node count, $N$ (nodes) & 16{,}000 / 32{,}000 / 100{,}000 & 0.446 / 0.223 / 0.071 \\
Storage horizon, $T$ (yr) & 5 / 10 / 15 & 0.446 / 0.223 / 0.149 \\
Avg block size, $B$ (MB) & 1.0 / 1.5 / 2.0 & 0.223 (invariant) \\
\bottomrule
\end{tabular}
\end{center}
\textbf{BTC price (the dominant omitted driver --- verified in independent implementation):} the ratio is linear in price. At the live N=32K baseline (capture price $\approx$ \$63{,}018 USD/BTC): \$30K$\rightarrow$0.106, \$45K$\rightarrow$0.159, \$62.9K$\rightarrow$0.222, \$100K$\rightarrow$0.354, \$200K$\rightarrow$0.707, \$500K$\rightarrow$1.768.

\textbf{Caveat --- price and fee level are not independent shocks.} The table above varies the BTC price holding the fee level (sat/vB) fixed, but in historical data fee levels and price co-move: bull markets raise both the price \emph{and} on-chain fee demand, so a price rise is typically accompanied by a fee-level rise, compounding the ratio's increase; a price fall in a demand crash carries fee levels down with it. The single-lever rows are \emph{ceteris-paribus} isolations for the model's arithmetic, not forecasts along a realized price path; the joint direction of price-and-fee co-movement is exactly the kind of dynamic question the v3.0 agenda (companion \texttt{future-directions-v3.md}) leaves open.

\textbf{Structural property:} SCCR is \textbf{invariant to block size} --- $c_b \propto 1/B$ while $L \propto B$, so $B$ cancels. Caveat: this invariance follows from attributing \emph{all} node cost per byte then multiplying back by bytes --- it reflects that the model is size-agnostic (contains no size-dependent economics), not that size is economically irrelevant.

\subsection{The knife-edge: node count and the strong claim}\label{sec:5.4}
The SCCR is homogeneous of degree 1 in its scale parameters: $\mathrm{SCCR} \propto (\mathrm{fee} \times \mathrm{price}) / (C \times T \times N)$. Two thresholds matter (derived and verified in the independent C implementation):
\begin{itemize}[nosep]
  \item The average SCCR crosses 1.0 at $N \approx$ 10{,}300 nodes (or BTC $\approx$ \$366{,}000) --- dated v2.0.0-era reference values (baseline 0.1719 at N=60K)
  \item The ``100\% of sampled blocks below $1\times$'' claim breaks at $N \approx$ 49{,}200 nodes (or BTC $\approx$ \$76{,}700) --- the highest-fee sampled block (13.75M sats, height 960469) sits at 0.82$\times$ coverage under the old N=60K
\end{itemize}
\textbf{Live recompute (2026-08-02 baseline, SCCR $=$ 0.2228 at N=32K):} the average now inverts at $N \approx$ 7{,}130 nodes or BTC $\approx$ \$283{,}000; the ``100\% below $1\times$'' break on the dated capture is unchanged at $N \approx$ 49{,}200. Both thresholds are reported; model-spec v2.0.1 retains the v2.0.0-era values (10.3K / \$366K).

\textbf{The node census --- a primary-source lower-bound census ($\geq$32{,}000 known addresses via Bitcoin Core \texttt{getnodeaddresses}), not a full enumeration.} We replaced the 60K assumption with data from a live Bitcoin Core node: a \textbf{\texttt{getnodeaddresses} RPC query} returned \textbf{32{,}000 known addresses} --- the RPC maximum, meaning the node's address manager is saturated at the cap and the true reachable set is \emph{at least} 32K. At census time the node also reported 8 live outbound P2P connections (2026-08-02), small relative to the 32K known-address lower bound. \textbf{We do not call this a complete census}: 32K is an artifact of the RPC cap, not a measured enumeration, and independent estimates span $\sim$10K--100K. The paper therefore reports the ratio as a \textbf{range across the true-N band}, with 32K as the best-available lower-bound estimate.

\textbf{Consequence of the real N=32{,}000 (recomputed, dated 156-block capture):}
\begin{center}
\begin{tabular}{lcc}
\toprule
Metric & At N=60K (old assumption) & At N=32K (primary-source lower-bound census) \\
\midrule
Average SCCR (dimensionless) & 0.156--0.172 & \textbf{$\sim$0.293} \\
Max per-block ratio (dimensionless) & 0.820 & \textbf{1.537} \\
Blocks below $1\times$ & 100.0\% & \textbf{98.7\%} (154/156) \\
\bottomrule
\end{tabular}
\end{center}
\textbf{This is a substantive finding, not a cosmetic one.} With a defensible node count, fees cover \emph{more} of the modeled storage cost than the 60K assumption implied ($\sim$29\% vs $\sim$17\%), and a handful of high-fee blocks now \emph{exceed} $1\times$ coverage. The direction of the headline is unchanged --- most blocks are still below $1\times$ --- but the strong form (``100\% below $1\times$'') does \textbf{not} survive the primary-source lower-bound census on the dated capture. \textbf{The defensible claim:} $\sim$22--29\% average coverage at N=32K, $\sim$98.7--100\% of sampled blocks below $1\times$, and $\sim$0.07--0.71 average across the true-N band (10K--100K). ($1\times$ remains a descriptive calibration point, not a normative target --- \S\ref{sec:1}.)

\textbf{Joint Monte Carlo on the headline} (\texttt{research/sccr\_monte\_carlo.py}, 10{,}000 samples, $N \sim \mathrm{Tri}(10K,150K,\text{mode }60K)$, $C \sim \mathrm{Tri}(\$500,\$2000,\text{mode }\$925)$, $T \sim \mathrm{Tri}(5,30,\text{mode }10)$, $P \sim \mathrm{Tri}(\$30K,\$120K,\text{mode }\$62.9K)$):
\begin{center}
\begin{tabular}{lc}
\toprule
Quantile & SCCR (dimensionless) \\
\midrule
P5 & 0.034 \\
P25 & 0.062 \\
P50 & 0.099 \\
P75 & 0.163 \\
P95 & 0.338 \\
Share below $1\times$ & \textbf{99.8\%} \\
\bottomrule
\end{tabular}
\end{center}
The median is below the deterministic point estimate because the node-count distribution's right tail (median N $\approx$ 86K) dominates. The result is robust in distribution.

\textbf{Current-N-band Monte Carlo} (\texttt{research/sccr\_monte\_carlo\_range.py}, 10{,}000 draws, $N \sim \mathrm{Tri}(10K,100K,\text{mode }32K)$ --- the paper's stated uncertainty band with the census as mode; $C \sim \mathrm{Tri}(\$600,\$1400,\text{mode }\$925)$, $T \sim \mathrm{Tri}(5,15,\text{mode }10)$, $P \sim \mathrm{Tri}(\$30K,\$120K,\text{mode }\$63{,}018)$; anchored at the live baseline and cross-checked against the frozen capture):
\begin{center}
\begin{tabular}{lc}
\toprule
Quantile & SCCR (dimensionless, live anchor) \\
\midrule
P5 & 0.071 \\
P25 & 0.117 \\
P50 & 0.173 \\
P75 & 0.260 \\
P95 & 0.467 \\
Share below $1\times$ & \textbf{99.9\%} \\
\bottomrule
\end{tabular}
\end{center}
The current-band median (0.17) sits near the N=32K point estimates (0.22--0.29), and the P5--P95 interval (0.07--0.47) brackets the deterministic N-band range ($\sim$0.07--0.71): the confidence interval and the range are two views of the same N-driven uncertainty. The old-N-band result above (99.8\% below $1\times$) remains as the historical conservative check.

\subsection{The Bandwidth Leg (v1 analytical bound)}\label{sec:5.5}
*(Added 2026-08-04.)* The storage leg prices \emph{stored} bytes; the bandwidth leg prices the \textbf{marginal propagation} of those bytes: every full node --- pruned or archival --- must download every byte of every block. Quantities live in \texttt{model-spec.json} v2.1.0 (\texttt{cost\_per\_gb} input; derived \texttt{bw\_GB\_yr}, \texttt{bw\_cost\_per\_year\_node}, \texttt{bw\_cost\_per\_year\_net}, \texttt{bw\_insc\_incr\_node}); the formula is \textbf{$B \times$ replication $\times$ \$/GB}:

\begin{center}
\begin{tabular}{lcl}
\toprule
Quantity & Value & Meaning \\
\midrule
\texttt{bw\_GB\_yr} = $B_{\mathrm{all,yr}}$/1e9 & \textbf{78.894 GB/yr} & full-chain bytes each node downloads \\
\texttt{cost\_per\_gb} (input) & \textbf{\$0.05/GB} & retail bandwidth proxy (AWS egress $\sim$\$0.09/GB, Hetzner $\sim$\$0.011/GB --- mid-range) \\
\texttt{bw\_cost\_per\_year\_node} & \textbf{\$3.94/yr per node} & marginal full-chain propagation, one node \\
\texttt{bw\_cost\_per\_year\_net} & \textbf{$\sim$\$126K/yr} & $\times$ $N=32$K lower-bound census \\
\texttt{bw\_insc\_incr\_node} & \textbf{\$0.024/yr per node} & inscription-incremental share (480 MB/yr $\times$ \$0.05/GB) --- reconciles the earlier ``$\sim$\$0.02/yr'' \\
\bottomrule
\end{tabular}
\end{center}

\textbf{Honest framing:} this is an \textbf{analytical bound, not a measurement.} Real node bandwidth bills are overwhelmingly flat-rate (residential/colocation unmetered), so \$/GB is a \emph{marginal economic proxy}, not a typical bill. The full-chain marginal cost is $\sim$\$3.94/yr per node ($\sim$\$126K/yr network at $\geq$32K nodes), and the inscription-incremental slice is $\sim$\$0.024/yr per node --- small in both frames, in the same direction as the pruning analysis verdict (\$0.02--0.05/yr unavoidable per node).

\subsection{The Validation Leg (v1 order-of-magnitude survey)}\label{sec:5.6}
*(Added 2026-08-04; full survey: \texttt{research/validation-cost.md}.)* Validation is the CPU work every node pays for every block since genesis --- PoW-header check (trivial), block rules, and \textbf{signature verification} (the dominant term) --- and it is the only leg strictly unavoidable at any replication scale (storage can be pruned; validation cannot be skipped). The v1 bound, \textbf{order-of-magnitude only} (band $\approx$ 0.5$\times$--5$\times$):

\begin{center}
\begin{tabular}{lcl}
\toprule
Term & Value & \\
\midrule
Blocks/yr ($R_{\mathrm{blocks}}$) & 52,596 & \\
Sig checks per block / throughput & $\sim$3--10K sigs/block vs $\sim$10--30K sigs/s (modern hw, libsecp256k1) & \\
Validation CPU per block & $\sim$0.1--1 s (cross-checked by initial-sync delta: 6--24 h / $\sim$1M blocks) & \\
Validation CPU per node per year & $\sim$1.5--15 h & \\
\textbf{Validation cost per node per year} & $\sim$\$0.5--\$5 (central $\approx$ \$1--2/yr at \$0.10--0.50/CPU-h) & \\
Network-wide ($N=32$K) & $\sim$\$16K--160K/yr & \\
\bottomrule
\end{tabular}
\end{center}

\textbf{Falsifiable claim (v1):} \emph{validation cost per full node per year is $<$\$100 --- bounded from above by the entire node budget $C=\$925$/yr, central estimate $\sim$\$1--2/yr.} Falsified by a measured benchmark showing $\geq \sim$200 h/yr steady-state validation CPU, or a hardware census showing validation is a \emph{binding} provisioning constraint. Literature anchored on Tschorsch \& Scheuermann (2016, IEEE COMST), Delgado-Segura et al.\ (2018, UTXO analysis), Bitcoin Core's \texttt{src/bench} suite. \textbf{Verdict:} validation is cheap per block, so cheap per node; its network total ($\sim$\$16K--160K/yr) is $\sim$3 orders of magnitude below the storage leg's modeled network burden --- which is itself the finding, not an assertion.

\subsection{The UTXO Leg (v1 measurement)}\label{sec:5.7}
*(Added 2026-08-04.)* \texttt{getblockstats} $\to$ \texttt{utxo\_size\_inc} --- the net byte delta of the UTXO set after each block --- was already captured by agent-06 but \textbf{dropped at DB write} (no \texttt{block\_stats} column). Fixed end-to-end: column added to \texttt{tools/db/schema.sql}, wired through \texttt{tools/db/init.js} (\texttt{insertBlockStats} + idempotent \texttt{ALTER TABLE} migration), live path in \texttt{tools/data-engineering/spool-consumer.js}, and a backfill (\texttt{tools/db/backfill-block-stats.js}) populated history from spool + raw captures. Measured v1 result (2026-08-04, \texttt{captured-data/bsahi.db}):

\begin{center}
\begin{tabular}{lcl}
\toprule
Metric & Value & \\
\midrule
\texttt{block\_stats} rows with \texttt{utxo\_size\_inc} & \textbf{165 unique heights} (backfill: 351 writes, 119 capture files) & \\
Avg net UTXO-set delta per block & $\sim$\textbf{29.9 KB/block} (range incl.\ genesis's 117 B artifact) & \\
Max height covered & 671,460 & \\
Live path & verified (handler test: height$\to$\texttt{utxo\_size\_inc} persisted) & \\
\bottomrule
\end{tabular}
\end{center}

\textbf{Status:} the leg is now \emph{measured} (the resource's size per block is in the queryable DB, growing live). What remains is the \emph{pricing} step --- attributing a \$/byte/yr cost to the UTXO delta --- which stays in \S\ref{sec:7} future work until then.

\section{Internal Validation, Correction, and Reconciliation (v2.0.0)}\label{sec:6}
Internal validation identified an inconsistency in the SCCR implementation, traced it to a \textbf{duplicated time-horizon term}, corrected the implementation, regenerated all reported values, and confirmed the qualitative conclusions unchanged. Four labeled steps: \textbf{Bug $\rightarrow$ Fix $\rightarrow$ Results $\rightarrow$ Reconciliation}.

\subsection{The bug}\label{sec:6.1}
\texttt{methodology.json} v1.0.0 and \texttt{storage-ratio.js} applied the horizon $T$ twice --- once dividing the denominator of \texttt{costPerBytePerYear} ($C / (B_{\mathrm{all,yr}} / T)$) and again in the lifetime-cost product ($\mathrm{bytes} \times c_b \times T$). This inflated modeled storage cost (USD/block) --- and deflated the ratio --- by exactly \textbf{$10\times$}. The error affected the magnitude, not the logical structure of the model.

\subsection{The fix}\label{sec:6.2}
$c_b$ is defined horizon-free ($C / \text{bytes-per-year}$); $T$ enters only through $L = c_b \times B \times T$. This is the canonical v2.0.0 definition, pinned in \texttt{research/model-spec.json}; no script may reintroduce $T$ into $c_b$.

\subsection{Results}\label{sec:6.3}
\begin{center}
\renewcommand{\arraystretch}{1.2}
\begin{tabular}{lcc}
\toprule
Quantity & Before (v1.0.0 formula, N=60K) & After (v2.0.0, N=60K) \\
\midrule
Cost per byte per year (USD/(byte$\cdot$yr)) & $1.17247\times10^{-7}$ & \textbf{$1.17247\times10^{-8}$} \\
Lifetime storage cost/node/block (USD/block) & \$1.759 & \textbf{\$0.176} \\
Network lifetime cost/block (USD/block) & \$105{,}522 & \textbf{\$10{,}552} \\
\textbf{Average SCCR (dimensionless)} & 0.0172 & \textbf{0.1719} \\
Blocks below $1\times$ & 100\% & \textbf{100\% (unchanged)} \\
\bottomrule
\end{tabular}
\end{center}
\textbf{Why the conclusion survived the correction.} Although the correction increased the estimated SCCR by $\sim10\times$ on a same-capture basis (and $\sim$11.5$\times$ across the as-reported headlines v1.0.0 \textbf{0.0149} $\rightarrow$ v2.0.0 \textbf{0.1719}), every sampled block still remained below the modeled full-cost threshold ($1\times$) under the then-assumed node count (N=60K). At the primary-source lower-bound census N=32K the average rose further to $\sim$0.225 on the 2026-08-02 live re-measure (168 blocks), with $\sim$99\% of sampled blocks still below $1\times$. The magnitude moved by an order of magnitude; the \emph{direction} of the finding did not.

\subsection{Reconciliation}\label{sec:6.4}
The two cost models (\texttt{storage-ratio.js} and \texttt{utxo\_cost\_model.py}) disagreed by $16.4\times$ (dimensionless). This decomposed as $164\times$ (dimensionless) denominator gap $\div$ $10\times$ (dimensionless) bug $=$ $16.4\times$. The models measure different quantities (block-average vs.\ marginal-inscription attribution); the gap is documented, not erased. See \texttt{verification\_appendix.md}.

\subsection{Reproducibility}\label{sec:6.5}
Every quantity in this paper is regenerated from \texttt{research/model-spec.json} (v2.0.1) by three independent implementations (JS, Python, standalone C); no script redefines a model constant, and the full capture log is retained in \texttt{captured-data/bsahi.db}. A frozen input capture and a cross-check script ship in \texttt{research/reproduce/} (see README).

\section{Limitations and Future Work}\label{sec:7}
\textbf{Limitations:}
\begin{enumerate}[nosep]
  \item The node count ($\geq$32K from the primary-source lower-bound census) is a \textbf{lower bound}, not a complete enumeration. SCCR is inversely proportional to N: at the live baseline the average inverts above $1\times$ only below $\sim$7.1K nodes, and the ``100\% below $1\times$'' claim breaks below $\sim$49K nodes at the dated capture (\S\ref{sec:5.4}).
  \item \textbf{Node costs are homogeneous --- and bundled.} Hardware/bandwidth/electricity vary by geography and operator; the model treats them as one scalar $C = \$925$/yr. Because $C$ is bundled ($C = C_{\mathrm{storage}} + C_{\mathrm{bandwidth}} + C_{\mathrm{misc}}$, \S\ref{sec:4.1}), the headline ratio is strictly a \textbf{storage-and-hosting coverage ratio}, and the bandwidth-vs-storage decomposition is a documented future refinement (RCIR and BCIR legs).
  \item 10-year horizon is an assumption; pruning shortens actual retention, permanent storage extends it.
  \item \textbf{No discounting.} A one-time fee (USD/block) is compared against an undiscounted 10-yr storage-cost sum; discounting the liability at $r=5\%$/yr ($8\%$/yr) reduces the present value by $\sim$27\% (45\%).
  \item Bandwidth is included in the fixed node cost ($C$) yet marginal bandwidth-propagation cost is excluded from the storage leg (fixed-vs-marginal distinction).
  \item Marginal vs.\ average attribution changes the per-byte cost by $164\times$ (dimensionless) --- the choice is explicit and documented.
  \item \textbf{Why storage at all?} The paper does not claim storage is Bitcoin's most important resource --- \textbf{it is simply the first measurable one}. Bandwidth, validation, UTXO, relay, and indexer serving are named research hypotheses with their own metrics (roadmap \S4), none yet measured.
\end{enumerate}
\textbf{Future work (the resource-pricing program; statuses updated 2026-08-04):} UTXO leg --- v1 DONE (measurement: \texttt{utxo\_size\_inc} persisted end-to-end, \S\ref{sec:5}.7; $\sim$29.9 KB/block avg, 165 backfilled heights); remaining: \$/byte pricing of the UTXO delta. Bandwidth leg --- v1 DONE (analytical bound: model-spec v2.1.0, \$3.94/yr per node marginal propagation, \S\ref{sec:5}.5); remaining: measured per-node bills. Validation leg --- v1 STARTED (OOM survey, \texttt{validation-cost.md}, <\$100/yr/node central \$1--2/yr, \S\ref{sec:5}.6); remaining: pinned Core benchmark + hardware census. Node distribution (\texttt{node\_geo}, 224 rows); BIP-110 pre/post measurement protocol if ever signaled; v3.0 economic-dynamics program (companion \texttt{future-directions-v3.md}, agenda summary \S\ref{sec:10}).

\subsection{What would falsify this framework?}\label{sec:7.1}
\textit{(Added 2026-08-03, post-advisor review --- see
\texttt{docs/decisions/2026-08-02-publication-decisions.md}.)} A framework that
cannot specify its own failure conditions is not a framework; it is a posture. We
state, in advance, the observations that would force us to retract or
substantially revise the claims in this paper. They operate at two levels --- the
\textbf{measurement} (SCCR, this paper) and the \textbf{framework} (the RIR
family and the ``no single resource market'' thesis; outline in
\texttt{research/framework-paper-outline.md}).

\textbf{Two external outcomes --- never conflated (added 2026-08-03,
third-reviewer refinement).} An external party's run of the reproduction
protocol can end in exactly two ways, and they are categorically different:
\begin{itemize}[nosep]
  \item \textbf{``Failed to reproduce''} --- an independent party following the
  published protocol from a clean clone gets a materially different number: not
  the published avg 0.2186, not the 0.07--0.71 band, or a per-block mismatch
  that cannot be reconciled. This is a \textbf{genuine falsification of the
  measurement}: the number is not trustworthy as stated, and the paper's
  headline falls until the discrepancy is reconciled. This is the only
  external-outcome class that blocks submission (falsifier 1 below; D5).
  \item \textbf{``Reproduced the number, disagrees with
  framing/assumptions''} --- an independent party runs the protocol, confirms
  the arithmetic (avg 0.2186, the band, the below-$1\times$ share), but
  disputes a modeling choice: the $C = \$925$/yr bundling, the $T=10$ horizon,
  storage-as-the-first-resource priority, or whether an unpriced-but-avoidable
  cost is an externality at all. This is \textbf{NOT a falsification}. It is
  honest scientific disagreement about documented assumptions --- the paper
  states these choices as assumptions (\S7, items 1--7) and anticipates exactly
  this class of challenge. It is folded into future revisions the way the Liu
  et al.\ (2021) prior work was reconciled (\S\ref{sec:8.2}): acknowledged,
  engaged, absorbed into the next version. The measurement stands.
\end{itemize}
The two outcomes must never be conflated. A framing disagreement is
\textbf{feedback, not a failed reproduction}: it is recorded in the
community-feedback triage (\texttt{research/community-review-plan.md} \S4
$\rightarrow$ \texttt{research/community-feedback.md}) and addressed in the
next revision, while the reproduced number stands. A failed reproduction is a
claim about the number itself and is handled under falsifier 1.

\textbf{Falsifiers of the SCCR measurement:}
\begin{enumerate}[nosep]
  \item \textbf{Independent implementations cannot reproduce the ratio.} The framework ships a frozen capture plus three independent implementations (JS/Python/C, verified agreeing per-block). If an uninvolved party running the protocol on the same inputs cannot reproduce SCCR $\approx$ 0.22 (or the discrepancy cannot be reconciled), the measurement is not trustworthy and the paper's headline falls. This is the external-reproduction critical path (\texttt{research/reproduce/external-reproduction.md}).
  \item \textbf{A better storage-cost model reverses the conclusion.} The denominator ($C \cdot T \cdot N / R_{\mathrm{blocks}}$) rests on a homogeneous node-cost model, an assumed $T=10$ horizon, and the $\geq$32K primary-source lower-bound census. If a defensible refinement --- measured pruned-vs-archival retention, regional cost heterogeneity, or corrected discounting --- moves the banded result ($\sim$22--29\% coverage, most blocks below $1\times$) to a qualitatively different conclusion, the paper's central claim is falsified in its current form.
  \item \textbf{Fees consistently exceed modeled long-term costs.} If sustained fee regimes (or re-measured 2017--2024 fee-peak years) show fees \emph{covering} modeled costs in a stationary way, the partial-internalization reading collapses into ``the market internalizes storage when it matters'' --- a different result, and the banded headline (\S\ref{sec:5}) would be wrong.
\end{enumerate}

\textbf{Falsifiers of the framework (the RIR family and the ``no single resource market'' thesis):}
\begin{enumerate}[nosep]
  \setcounter{enumi}{3}
  \item \textbf{Resource-cost attribution is shown economically inappropriate.} If the planned attribute-pricing regression (companion \texttt{future-directions-v3.md} \S2 Q2) shows the single fee price carries \emph{no} persistence signal, then per-resource internalization ratios are category errors, not measurements --- the RIR family premise fails even though SCCR-as-computed may stand.
  \item \textbf{The pruned-vs-archival census shows the storage burden is avoidable at scale.} The $T=10$ replicated-storage denominator assumes archival retention across the network. A measured pruning distribution showing most nodes avoid most of the burden (companion note \texttt{archival-vs-pruned-note.md}) would reframe SCCR as an upper bound on an avoidable cost --- a genuine falsification of the externality reading, though not of the measurement.
  \item \textbf{Measured response functions close the dynamic loop at or above $1\times$.} If node-entry/exit and fee-demand response functions (companion \texttt{future-directions-v3.md} \S2 Q1) are measured and exhibit a stable fixed point at or above $1\times$, the persistent-partial-internalization equilibrium hypothesis is falsified.
\end{enumerate}

We do not believe any of these falsifiers currently obtain; we list them so the
reader can check, and so the framework is never mistaken for an unfalsifiable
claim.

\subsection{Assumption taxonomy --- deliberate choices (Type C) vs.\ empirical risks (Type M)}\label{sec:7.2}
\textit{(Added 2026-08-03, fourth-reviewer deliverable.)}
\S\ref{sec:7.1} draws a categorical line between two external outcomes --- a
\textbf{failed reproduction} (the number is wrong) and a \textbf{reproduced
number with disagreed framing} (the number stands). That line is only
enforceable if the paper's own assumptions are pre-sorted into the same two
buckets \textbf{before} the challenge arrives --- not argued case-by-case under
time pressure when an external review lands. Every material assumption in this
paper therefore belongs to exactly one of two categories, stated here in advance:

\textbf{The crisp test.} A challenge is classified by what it does to the
assumption's current value, not by who makes it:
\begin{itemize}[nosep]
  \item \textbf{If a critic simply asserts a different value or a different
  modeling choice} (T = 20, marginal attribution, ``validation matters more'')
  --- there is no evidence the current value is wrong. That is \textbf{Type C
  territory: recorded as feedback} (community-feedback triage,
  \texttt{research/community-review-plan.md} \S4), engaged in the next revision,
  and \textbf{the measurement stands}.
  \item \textbf{If the critic shows the current value or measurement is wrong}
  --- with evidence, a reproduction mismatch, or better data (a complete node
  census, a corrected fee capture, a measured cost decomposition) --- that is
  \textbf{Type M territory: a falsification candidate} under \S\ref{sec:7.1}
  falsifiers 1--3, and \textbf{the headline is at risk} until the discrepancy is
  reconciled.
\end{itemize}

\textbf{Type C --- ``Assumption I chose'' (deliberate modeling choice).} A
selection among defensible alternatives, made explicit and documented where it
appears (\S7 items 1--7). Not falsifiable by asserting a different choice; a
reviewer preferring another selection records \textbf{feedback}, not
falsification. Reclassification to Type M requires showing the choice is
internally inconsistent with the paper's own stated method, or rests on an
empirical premise that is wrong.

\textbf{Type M --- ``Assumption that might be mismeasured'' (empirical risk).} An
empirical quantity, measurement, or bound embedded in the model. Falsifiable: if
the measurement is wrong, the conclusion may shift. A critic who shows the
current value/measurement is wrong opens a \textbf{falsification candidate} ---
handled under \S\ref{sec:7.1} falsifiers 1--3; the headline falls until
reconciled.

\begin{table}[h]
\centering
\small
\renewcommand{\arraystretch}{1.3}
\begin{tabular}{@{}>{\raggedright\arraybackslash}p{0.7cm}>{\raggedright\arraybackslash}p{2.9cm}>{\raggedright\arraybackslash}p{2.0cm}>{\raggedright\arraybackslash}p{5.3cm}>{\raggedright\arraybackslash}p{5.0cm}@{}}
\toprule
\textbf{\#} & \textbf{Assumption} & \textbf{Type} & \textbf{Why it's that type} & \textbf{What would reclassify or falsify it}\\
\midrule
1 & Storage horizon $T = 10$ yr & \textbf{C} & A selection within a defensible range (pruning shortens retention, permanent storage extends it, \S7 item 3); sensitivity disclosed (\S\ref{sec:5.3}: $T=5/10/15 \rightarrow 0.446/0.223/0.149$), direction robust & A measured network-wide retention distribution showing effective retention is an order of magnitude shorter --- reframes the externality reading (falsifier 5); does not make $T=10$ an error\\
2 & Replication factor $N = 32$K & \textbf{M} & An empirical bound from the primary-source census (addrman cap at 32{,}000 known addresses), explicitly \emph{not} a complete enumeration (\S\ref{sec:5.4}); true set estimated 10K--100K; SCCR $\propto 1/N$ & A defensible complete census showing the true reachable set differs materially from 32K in a direction that moves the $\sim$0.07--0.71 band; the ``$\sim$99--100\% below $1\times$'' claim already breaks at $N \approx 49$K (\S\ref{sec:5.4})\\
3 & \textbf{C bundling} ($C = C_{\mathrm{storage}} + C_{\mathrm{bandwidth}} + C_{\mathrm{misc}}$) & \textbf{M (leaning --- the blurry case)} & Bundling is presented as a simplification (Type C-like), but it embeds an implicit empirical claim that the components are as measured --- and the component sum is bandwidth-dominated (600 + 166.67 + 157.68, \S\ref{sec:4.1}), so the \emph{storage} reading of the headline is fragile if the decomposition is wrong & A measured node-cost decomposition (operator surveys, hosting bills) showing the storage share is materially different, or a double-counted component --- reframes the headline from ``storage-and-hosting coverage'' (\S\ref{sec:4.1}) and re-bands the ratio\\
4 & Average block size $B_{\mathrm{block}} = 1.5$MB & \textbf{M (narrow surface)} & An empirical quantity from the fee-history capture --- but \S\ref{sec:5.3} shows SCCR is invariant to $B$ ($B$ cancels), so a wrong value changes the per-byte presentation ($c_b$), not the headline ratio & A material mis-statement of the capture (wrong byte basis in the fee history) propagating through $c_b(t)$; a data-quality flag, not headline-moving under the current spec\\
5 & \textbf{``Fees paid'' definition} ($\mathrm{fee\_USD} = \mathrm{avgFees} \times \mathrm{USD/BTC}$; subsidy excluded, \S\ref{sec:4.1}) & \textbf{M} & The numerator is a measurement from the live capture (mempool.space 24h block-fee history); the definition is explicit but the value is empirical & A reproduction mismatch on the fee side, a corrected capture, or evidence the conversion/attribution is wrong --- direct falsifier 1/3 territory: the number falls until reconciled\\
6 & \textbf{Storage-first sequencing} & \textbf{C} & A program ordering, explicitly ``the first measurable one,'' not ``the most important'' (\S7 item 7) & Only evidence that storage is \emph{not} reproducibly measurable (falsifier 1) touches it; asserting validation/UTXO is more important is feedback, not falsification\\
7 & \textbf{Average-vs-marginal attribution} & \textbf{C} & A documented choice (the $164\times$ denominator gap is disclosed; both branches reported, \S\ref{sec:4.2}); average answers the accounting question, marginal the optimization question & Showing the marginal object is the only economically correct one for the externality claim --- the paper reports both, so this is engagement, not error\\
8 & \textbf{Homogeneous node cost structure} & \textbf{C (with an empirical seam)} & The single-scalar structure is a deliberate simplification, documented (\S7 item 2); heterogeneity is a refinement direction, not an error & The \emph{value inside} the structure ($C = \$925$/yr) is empirical --- a defensible regional refinement that moves the band is already named falsifier 2; the structural choice itself is not falsifiable by asserting heterogeneity\\
9 & \textbf{No discounting; constant cost} & \textbf{C} & A documented choice with disclosed sensitivity (\S7 item 4): discounting at $r=5\%/8\%$ cuts the liability's present value by $\sim$27\%/45\%, moving the average up ($\sim$0.30--0.53 at $N=32$K) --- still below the $1\times$ average threshold; the majority-below-$1\times$ direction is robust & A corrected treatment that flips the headline's direction --- none identified under the disclosed sensitivity; a critic asserting discounting records feedback, not falsification\\
\bottomrule
\end{tabular}
\end{table}

\textbf{The blurry case, named honestly.} Row 3 (C bundling) is the one
assumption that sits on the line. The \emph{decision to bundle} is a choice
(Type C); the \emph{claim that the bundle's components are as measured} is
empirical (Type M). Because the bandwidth leg (600) is the largest component, the
storage-specific reading of the headline depends on a decomposition that is
currently \emph{assumed}, not measured. We therefore classify C bundling
\textbf{leaning Type M} --- the safe default: an external challenge to the
decomposition gets the full falsification response, not the feedback bucket.

\textbf{Why the taxonomy is stated in advance.} Every assumption-level critique
this paper receives will be routed by this table before it is answered: Type C
critiques are engaged as feedback in the next revision (the reproduced number
stands); Type M critiques are treated as falsification candidates until the
evidence is reconciled (\S\ref{sec:7.1} falsifiers 1--3). This is the mechanism
that keeps ``reproduced the number, disagrees with framing'' from silently
becoming ``the paper might just be wrong'' --- and keeps a genuine mismeasurement
from being dismissed as a mere framing disagreement. The plan-of-record mirror
lives in \texttt{roadmap.md} \S12.


\section{Economics and Related Work}\label{sec:8}
\subsection{Is this an externality?}\label{sec:8.1}
Following standard environmental-economics usage (Pigou, 1920; Coase, 1960), a \textbf{negative externality} arises when a transaction's cost is borne by parties who did not consent. Applied to Bitcoin: an inscription's fee is paid by its creator; the long-term storage cost is borne by node operators who neither created the transaction nor were compensated for it. Two qualifications are stated honestly:
\begin{enumerate}[nosep]
  \item \textbf{Voluntary participation.} Node operators choose to run nodes (and may prune). The Pigouvian case is therefore weaker than for a physical externality --- the correct framing is an \emph{unpriced but avoidable} cost, not an imposed one. \textbf{Voluntary participation weakens the welfare interpretation, but not the measurement} --- the ratio still quantifies what the fee market covers of a real, recurring cost borne by the storage-bearing class, whatever the operator's freedom to exit.
  \item \textbf{Measurement, not pricing.} This paper \emph{measures} a ratio; it contains no price mechanism and proposes no policy.
\end{enumerate}
\subsection{Related work and novelty}\label{sec:8.2}
\begin{itemize}[nosep]
  \item \textbf{Aronoff, Praizner, Sabouri (2026), arXiv:2604.17183} --- structural VCG fee model; treats the mempool as a market for scarce block space; does not model storage externalities.
  \item \textbf{Liu, Fang, Cheung, Cai, Huang (2021), arXiv:2103.05866} --- \emph{``An Incentive Mechanism for Sustainable Blockchain Storage.''} Argues storage costs have ``in general not been properly compensated by the users' transaction fees.'' \textbf{This is the closest prior work} and we acknowledge it directly: our contribution is not the observation (already argued in 2021), but a \emph{measured, reproducible, Bitcoin-live-data quantification} (the SCCR) with a reconciliation of cost models and a knife-edge sensitivity bound. \textbf{We do not claim the observation is novel; we claim the measurement is.}
  \item \textbf{Ethereum state-rent / gas-as-state-pricing} and \textbf{Sompolinsky--Zohar} are adjacent threads we build toward but do not model.
\end{itemize}
\subsection{The efficient-markets objection (named rebuttal)}\label{sec:8.3}
\textbf{Objection:} if block space is priced at the margin by a market that clears, then marginal cost is internalized by definition --- the fee IS the price, so there is no externality to measure.

\textbf{Rebuttal:} this conflates two different horizons. The fee market prices \emph{inclusion in the next block} --- a static, congestion-clearing price with a $\sim$10-min horizon. The storage cost the SCCR measures is a \emph{recurring} cost over an indefinite horizon (the data persists in every full node's history). A market can clear for the short-horizon good (block space now) while not pricing the long-horizon good (permanent replicated storage). The paper's measurement question is precisely whether the one-time congestion price also covers the recurring storage cost; \S\ref{sec:5} shows the empirical answer (banded $\sim$22--29\% coverage at N=32K). The efficient-markets framing is therefore not contradicted --- it is \emph{scoped}: it describes the short-horizon market, and the paper measures the long-horizon gap.

\textbf{Designer-intent corroboration (primary sources).} The scoped reading is not a reinterpretation --- it matches the system's own design intent. The whitepaper defines the fee solely as an incentive to create blocks: \emph{``If the output value of a transaction is less than its input value, the difference is a transaction fee that is added to the incentive value of the block''} (\S6), and storage is explicitly designed \emph{around}, not priced: the 80-byte block-header estimate (\emph{``80 bytes $*$ 6 $*$ 24 $*$ 365 $=$ 4.2MB per year''}) is dismissed with \emph{``storage should not be a problem''} (\S7) --- a claim about block headers kept in memory for SPV, not full-chain replication, which is precisely the storage this paper measures. When Satoshi later moved to protect disk space, he used quantity control, not price: the block-size threshold should be kept \emph{``lower as a circuit breaker''}, explicitly to \emph{``limit the amount of wasted disk space in that event''} (BitcoinTalk post 441, Sep 2010), and the cap was designed to be raised by a block-height trigger, \emph{``if (blocknumber $>$ 115000) maxblocksize $=$ largerlimit''} (post 485, Oct 2010). The efficient-markets objection therefore fails on the designer's own terms: the single price was specified as an inclusion incentive; storage was handled --- where it was handled at all --- by quantity controls, never by price. (Full verification, exact sources, the apocryphal-quote flag, and a falsifiable-claims table derived from these primary texts: companion note \texttt{research/satoshi-primary-source-note.md}; the ``node equilibrium self-corrects'' claim maps to roadmap Q1.)

\subsection{Terminology}\label{sec:8.4}
This paper's primary name for the metric is \textbf{Storage Cost Coverage Ratio (SCCR)} --- a fee-to-modeled-cost comparison, \emph{not} a coverage ratio in the insurance/solvency sense. The broader economics family generalizes it: \textbf{Cost Internalization Ratio} is the family name for ``the share of a long-lived resource cost internalized by the one-time fee,'' of which SCCR (storage) is the first measured member (Metric \#1 of the RIR family; companion \texttt{future-directions-v3.md} \S2 Q3). We use SCCR / Storage Cost Coverage Ratio as the established name throughout; Cost Internalization Ratio appears here once to locate the family.

\section{Conclusion}\label{sec:9}
Whether the fee market internalizes long-lived resource costs is an \textbf{open empirical question}, and this paper contributes the first reproducible measurement using a \textbf{primary-source lower-bound census} ($\geq$32{,}000 known addresses via Bitcoin Core \texttt{getnodeaddresses}): a storage-cost-internalization ratio of \textbf{$\sim$0.22--0.29 at N=32K} (167--156 blocks, dated captures), with \textbf{$\sim$98.7--100\% of sampled blocks below the $1\times$ threshold} at that node count. Because the replication factor $N$ is the dominant input and is only bounded ($\geq$32K; independent estimates 10K--100K), the honest headline is a \textbf{range}: fees currently cover roughly \textbf{7\% to 71\%} of the modeled 10-year replicated storage cost depending on the true node count, with the \emph{direction} (most blocks below $1\times$ under any $N \geq \sim$32K census) robust to every correction made in this paper's review. \textbf{Status:} external reproduction is pending (D5, the only submission blocker); every figure is stated to the precision the current evidence licenses, and the N-band range carries the dominant uncertainty.

\section{Future Directions (v3.0) --- the agenda, in a companion}\label{sec:10}
The paper's measurement question is answered in \S\ref{sec:5}; its falsifiers are \S\ref{sec:7.1}. The program's forward agenda --- the v3.0 economic-dynamics questions, their first answers, and the cross-chain generalization --- \textbf{deliberately lives outside this core paper} in the companion \texttt{research/future-directions-v3.md} (plan-of-record: \texttt{roadmap.md} \S8/\S9/\S11). This keeps the paper a \emph{measurement}, not a research-program pitch. The eight v3.0 questions (headline model output at the live baseline, SCCR $=$ 0.2228 at N=32K, T=10, C=\$925/yr):
\begin{center}
\renewcommand{\arraystretch}{1.2}
\begin{tabular}{cll}
\toprule
\# & Question & Headline model output (live baseline) \\
\midrule
1 & Storage horizon: SCCR at $T = 5, 10, 20, 30, 50$ yr & 0.446 / 0.223 / 0.111 / 0.074 / 0.045 (inverse-linear) \\
2 & Storage $10\times$ cheaper: $C = \$92.5$/yr & SCCR $=$ 2.228 --- the gap flips sign \\
3 & BTC $= \$500{,}000$ & SCCR $=$ 1.768; average crosses $1\times$ at $\sim$\$283K \\
4 & Fees sustained at 100 sat/vB for 5 yr & SCCR $\approx$ 11.2 (22.4 at $T{=}5$); crosses $1\times$ at $\sim$9 sat/vB \\
5 & Lightning moves 90\% of payments off-chain & Two-sided; net effect open \\
6 & Nodes double / triple: $N = 32K \to 64K \to 128K$ & 0.223 $\to$ 0.111 $\to$ 0.056 (inverse-linear) \\
7 & Can SCCR reach 1 without protocol changes? & Historically yes: above $1\times$ in 2017--2024 fee-peak years \\
8 & What is the equilibrium? & \textbf{Open.} The model measures a ratio; it does not close the loop \\
\bottomrule
\end{tabular}
\end{center}
\textbf{Open-question honesty.} Nothing in this section claims the fee market \emph{will} internalize these costs, nor that it \emph{should}. The claim is narrower: the question is now well-posed, the drivers are identified, and each driver's lever on the ratio is measured.

\begin{thebibliography}{9}
\bibitem{pigou1920} Pigou, A.~C. \emph{The Economics of Welfare}. Macmillan, 1920.
\bibitem{coase1960} Coase, R.~H. ``The Problem of Social Cost.'' \emph{Journal of Law and Economics} 3 (1960): 1--44.
\bibitem{aronoff2026} Aronoff, D., Praizner, J., Sabouri, S. ``A Model and Estimation of the Bitcoin Transaction Fee.'' arXiv:2604.17183, 2026.
\bibitem{liu2021} Liu, Y., Fang, Z., Cheung, M.~H., Cai, W., Huang, J. ``An Incentive Mechanism for Sustainable Blockchain Storage.'' arXiv:2103.05866, 2021.
\bibitem{bip141} BIP~141: Segregated Witness. Bitcoin Improvement Proposal, 2015--2017.
\bibitem{sompolinsky2015} Sompolinsky, Y., Zohar, A. ``Secure High-Rate Transaction Processing in Bitcoin.'' \emph{FC 2015}.
\bibitem{nakamoto2008} Nakamoto, S. ``Bitcoin: A Peer-to-Peer Electronic Cash System.'' 2008. \texttt{https://bitcoin.org/bitcoin.pdf} (\S6 fee-as-incentive; \S7 ``storage should not be a problem'').
\bibitem{nakamoto_emails} Nakamoto, S. Emails to the Cryptography Mailing List, Oct--Nov 2008. Satoshi Nakamoto Institute archive: \texttt{https://satoshi.nakamotoinstitute.org/emails/cryptography/} (esp.\ \#2 bandwidth-in-dollars, \#13 fees-as-inclusion-incentive).
\bibitem{nakamoto_posts} Nakamoto, S. BitcoinTalk posts 188, 287, 441, 485 (2010). Satoshi Nakamoto Institute archive: \texttt{https://satoshi.nakamotoinstitute.org/posts/} (post 188: ``never more than 100K nodes'' + node equilibrium; post 441: block-size threshold as ``circuit breaker''; post 485: phased block-size increase).
\end{thebibliography}

\vfill
\noindent\footnotesize\textit{Bitcoin Sahi Research Council --- The Bitcoin Block Space Problem (Working Paper v2.2.0). Component docs: problem\_statement \textbullet{} bip141\_analysis \textbullet{} pruning\_externality\_analysis \textbullet{} utxo\_cost\_note \textbullet{} verification\_appendix \textbullet{} history-of-bitcoin. Conversion note: see header comment for conversion status.}

\end{document}
